\section{Method}\label{sec:method}

The article-track pipeline has six stages. Each stage corresponds to a
self-contained module in the companion repository (the module name is
indicated parenthetically below).

\subsection{Stage 1 -- Design of Experiments}
\label{sec:method-doe}
A user-supplied or library-generated experimental design
(\texttt{doe\_xgb.design}) covers the hyperparameter space. The
canonical artifact is the Minitab face-centered central composite
design ($\CCDFC$) of the dissertation, with $k{=}7$ XGBoost
hyperparameters, $n{=}88$ runs and four center points; the
\texttt{DesignProvider} also generates fractional factorial,
Box--Behnken, Latin hypercube, Sobol, and simplex-lattice designs
in pure NumPy. Each artifact carries matrix diagnostics (rank,
condition number, coverage of the coded box) and a SHA-256 of the
source CSV so reviewers can replay it byte-for-byte.

\subsection{Stage 2 -- Cross-validated evaluation}
\label{sec:method-cv}
Each design row is evaluated under stratified $k$-fold CV
(default $k{=}5$); per-fold metrics are persisted in a long-format
CSV. Quality metrics include accuracy, precision, recall, specificity,
F1, ROC-AUC, PR-AUC, balanced accuracy, MCC, Brier score, and
expected calibration error (ECE). Wall-clock cost is the mean
fit-plus-predict time per fold. The single source of truth for
hyperparameter naming and types is \texttt{doe\_xgb.config}.

\subsection{Stage 3 -- Factor model with Varimax rotation}
\label{sec:method-fa}
The article-track factor model
(\texttt{doe\_xgb.factor\_model}) supports four modes:
\textsc{auto} (Kaiser eigenvalue $>1$ \emph{or} cumulative variance
threshold), \textsc{fixed} (explicit number of factors -- the
article default for the headline run is three),
\textsc{manual} (user-declared constructs such as
Quality / Cost / Robustness), and \textsc{none}
(no FA; raw metrics flow through unchanged). All modes return
loadings, factor scores, eigenvalues, explained variance,
communalities, the rotation matrix, the construct map, and a KMO
estimate where defined.

\subsection{Stage 4 -- Objective specification}
\label{sec:method-objectives}
Every objective entering the optimization is declared through an
\texttt{ObjectiveSpec} (\texttt{doe\_xgb.objectives}) with required
fields \texttt{direction}
$\in \{\textsc{minimize},\textsc{maximize},\textsc{target}\}$,
\texttt{transform}
$\in \{\textsc{raw},\textsc{zscore},\textsc{log1p},\textsc{factor\_score},\textsc{fmse},\textsc{custom}\}$,
\texttt{role}
$\in \{\textsc{primary\_nbi},\textsc{post\_optimization},\textsc{reporting},\textsc{constraint}\}$,
and optional \texttt{target}, \texttt{weight}, \texttt{group},
\texttt{bounds}, and \texttt{normalization}.
A canonicalizer maps every spec to a minimization callable. For
\textsc{target} / FMSE objectives we use Eq.~21 of
\citet{pereira2025eaai}:
\begin{equation}
  f^{\text{FMSE}}_i(\bm x) \;=\; \bigl(\hat F_i(\bm x) - T_i\bigr)^{2}
                                  \;+\;\sigma^{2}_{\hat F_i},
  \label{eq:fmse}
\end{equation}
where $T_i$ is the target value (e.g., the per-objective surrogate
optimum) and $\sigma^{2}_{\hat F_i}$ is the factor-variance term
implied by the loadings $\bm L$ of stage 3. The article-track
default for $\VRF$ objectives is \textsc{target} + FMSE; raw ML
metrics (e.g., \emph{Time\_MeanFold}, \emph{Accuracy\_Mean}) use
\textsc{raw} with explicit \textsc{minimize} / \textsc{maximize}.

\subsection{Stage 5 -- Surrogate fitting}
\label{sec:method-surrogate}
For process variables (CCDFC, Box--Behnken, factorial), we fit a
full quadratic RSM with strong-hierarchy backward elimination at
$\alpha{=}0.05$ (\texttt{ProcessQuadraticRSM}). For mixture variables
on the simplex (used by the post-optimization stage of
Section~\ref{sec:method-mbpa}), we fit a Scheff\'{e} canonical
mixture polynomial of order linear / quadratic / special cubic / cubic
(\texttt{MixtureScheffeModel}). Backward elimination is intentionally
disabled for Scheff\'{e} models because the canonical interpretation
depends on the full term set. A
\texttt{validate\_for\_model} routine warns or rejects incompatible
design/family combinations (mixture model on CCDFC; process quadratic
on simplex; quadratic on Latin hypercube).

\subsection{Stage 6 -- True $N$-objective NBI}
\label{sec:method-nbi}
Let $\hat{\bm F}(\bm x) = (\hat F_1(\bm x),\dots,\hat F_q(\bm x))$ be
the canonicalized minimization surrogates. We compute per-objective
anchors $\bm x^{*}_j = \arg\min_{\bm x\in\Omega} \hat F_j(\bm x)$ and
evaluate every objective at every anchor to obtain the payoff matrix
$\bm{\Phi}^{\text{full}}_{ij} = \hat F_i(\bm x^{*}_j)$. The utopia
point is its diagonal,
$\bm F^{U}_i = \hat F_i(\bm x^{*}_i)$, and the CHIM matrix is
$\bm{\Phi} = \bm{\Phi}^{\text{full}} - \bm F^{U}\bm{1}^{\top}$. The
quasi-normal direction defaults to the Das--Dennis convention,
\begin{equation}
  \hat{\bm n} \;=\; \dfrac{-\bm{\Phi}\,\bm{1}}{\|\bm{\Phi}\,\bm{1}\|},
  \label{eq:nhat}
\end{equation}
with a Gram--Schmidt orthogonal-complement variant available as a
diagnostic. For a weight vector $\bm\beta$ on the simplex the
sub-problem is
\begin{equation}
  \begin{aligned}
    \max_{\bm x \in \Omega,\, t \ge 0} \quad & t \\
    \text{s.t.}\quad & \hat{\bm F}(\bm x) - \bm F^{U}
       - \bm{\Phi}\bm\beta - t\,\hat{\bm n} = \bm 0,
  \end{aligned}
  \label{eq:nbi-subproblem}
\end{equation}
solved as $q$ scalar equality constraints and $\bm x$ box bounds via
SLSQP with multistart. Weight vectors come from a Scheff\'{e}
simplex-lattice $\{q,m\}$ (default $m{=}50$ for $q{=}2$, $m{=}10$ for
$q{=}3$, giving 51 and 66 sub-problems respectively). The residual
$\|\hat{\bm F}(\bm x^{*}(\bm\beta)) - \bm F^{U} - \bm{\Phi}\bm\beta -
t\,\hat{\bm n}\|$ is reported for every sub-problem so reviewers can
audit constraint satisfaction.

\paragraph{Methodology pin.} The $N$-objective generality is
verified by automated tests for $q\in\{2,3,4,5\}$ on synthetic
quadratics with closed-form anchors; residual norms remain below
$10^{-3}$ across the entire test matrix.

\subsection{Stage 7 -- Conditional MBPA post-optimization}
\label{sec:method-mbpa}
A flat or compressed primary frontier reduces the practical value of
$\NBI$. We compute six diagnostics on the front -- normalized
objective range, average pairwise distance, count of unique
non-dominated points, weight concentration, curvature score, spread
$\Delta$ -- and trigger MBPA when any one fails its threshold
(\texttt{doe\_xgb.diagnostics}). When triggered
(\texttt{doe\_xgb.post\_optimization}) we fit Scheff\'{e} canonical
mixture surrogates over the simplex of weights for the
generalized distance to utopia
\citep[Eq.~29]{pereira2025eaai} and the Shannon entropy
$S(\bm w) = -\sum_i w_i \ln w_i$, then minimize $\mathrm{GD}(\bm w) -
S(\bm w)$ inside an elliptical interior constraint that excludes
vertex weights. Default radii are $0.30$, matching the EAAI
recommendation. When MBPA is skipped we still write
\texttt{frontier\_quality.json} so the decision is auditable.

\subsection{Final selection}
\label{sec:method-selection}
A pluggable selection rule (\texttt{doe\_xgb.selection}) chooses the
final candidate. The article-track default is
\textsc{distance\_to\_utopia}; \textsc{knee},
\textsc{utility}, \textsc{lexicographic}, and \textsc{topsis} are
available. Benchmarks retain \textsc{max\_accuracy} so the
dissertation tables remain comparable. The full Pareto front is
always persisted regardless of the selection rule.

\subsection{Legacy weighted-sum ablation}
\label{sec:method-legacy}
The dissertation-era weighted-sum scalarization
(\texttt{doe\_xgb.scalarization.run\_nbi\_weighted\_sum}) is preserved
verbatim and exposed as an ablation baseline. It is never referred to
as NBI in the article text or in the article-track code; the
discrepancy is documented in
\texttt{docs/METHODOLOGY\_DECISIONS.md} D1.
